\documentclass[12pt, a4paper]{article}
\usepackage[french]{babel}
\usepackage[utf8]{inputenc}
\usepackage[T1]{fontenc}
\usepackage{geometry}

\usepackage{hyperref}
\usepackage{csquotes}
\usepackage[style=verbose, backend=biber]{biblatex}
\addbibresource{biblio.bib}

\geometry{margin=1in}

\begin{document}

\pagestyle{empty}

\begin{center}
    \large \textbf{Résumé : Self-Supervised Learning from Images with a Joint-Embedding Predictive Architecture} \\
    \vspace{0.5cm}
    Cédric GAUTHERET, Romain AMIGON et Noé BARRANGER
\end{center}

I-JEPA \footcite{ijepa23} est une nouvelle architecture d'apprentissage auto-supervisé conçue pour apprendre des représentations d'images hautement sémantiques. Contrairement aux approches classiques, ce modèle ne dépend pas d'augmentations de données complexes (SimCLR, DINO, etc.) i.e. méthodes contrastives ni de méthodes génératives (GAN ou MAE) basées sur la reconstruction de pixels. Les auteurs visent à démontrer qu'en prédisant les parties manquantes d'une image directement dans un espace de représentation latent, le système peut capturer l'essence sémantique d'une scène tout en éliminant les détails de bas niveau inutiles, améliorant ainsi la généralisation du modèle et accélérant sont entraînement.

Le I-JEPA repose sur trois composants clés : un \textbf{context-encoder}, un \textbf{target-encoder} et un \textbf{predictor} chacun étant basé sur un modèle de ViT (le predicteur étant plus léger). Le coeur de l'innovation est le \textbf{masquage multi-blocs}, où le modèle doit prédire les représentations de plusieurs gros blocs cibles continus à partir d'un unique bloc de contexte informatif au lieu de petits blocs disjoints comme c'est le cas pour le MAE. Pour éviter le phénomène d'effondrement des représentations, les poids de l'encodeur cible sont mis à jour via \textbf{EMA} basée sur les poids de l'encodeur de contexte, sans calcul de gradient direct. La perte est calculée comme la distance $L^2$ moyenne entre les représentations prédites et les représentations cibles dans l'espace latent. Pour mettre à jour les poids du context-encoder et predictor, on utilise une descente de gradient reposant sur AdamW avec un \textit{CosineAnnealing scheduler} qui permettent de moduler intelligemment le taux d'apprentissage. 

Pour le pré-entraînement non supervisé, les auteurs utilisent l'ensemble de données ImageNet-1K. La robustesse et la polyvalence des représentations apprises sont ensuite évaluées sur une grande variété de tâches. Cela inclut des tests de classification sur des jeux de données tels que CIFAR100 ou Places205.

Les résultats expérimentaux montrent que I-JEPA est non seulement plus performant que les méthodes de reconstruction de pixels (comme MAE) pour la classification d'images, mais aussi beaucoup plus efficace en termes de ressources. Par exemple, l'entraînement d'un modèle ViT-Huge/14 nécessite moins de 1200 heures GPU, soit une efficacité décuplée par rapport à MAE. En outre, I-JEPA surpasse nettement les méthodes basées sur l'invariance sur les tâches de bas niveau nécessitant une précision géométrique, comme le comptage d'objets et la prédiction de profondeur. Les visualisations montrent que le prédicteur capture correctement les structures globales et les poses des objets.

L'une des forces majeures de I-JEPA réside dans sa capacité à apprendre des caractéristiques sémantiques riches de manière efficace sans introduire les biais induits par les augmentations de données (flou et bruit). Sa polyvalence lui permet de briller tant sur des tâches sémantiques que sur des tâches locales. Cependant, le modèle présente certaines limites, notamment la sensibilité de la performance liée à l'architecture légère du prédicteur. De plus, bien que le pré-entraînement soit globalement plus rapide, chaque itération d'entraînement comporte un léger surcoût lié au calcul des cibles dans l'espace latent comparativement aux méthodes prédisant directement les pixels.


\newpage

\begin{center}
    \large \textbf{Résumé : Masked Autoencoders Are Scalable Vision Learners} \\
    \vspace{0.5cm}
    Cédric GAUTHERET, Romain AMIGON et Noé BARRANGER
\end{center}

Les auteurs dans \footcite{mae22} s'attaquent à la problématique du pré-entraînement auto-supervisé en vision par ordinateur, motivés par l'écart de performance avec le NLP, où des méthodes comme BERT ont démontré la puissance du masked autoencoding. Pour ce faire, ils proposent une approche nommée Masked Autoencoder (MAE), reposant sur deux choix de conception centraux. Premièrement, une architecture encodeur-décodeur asymétrique : l'encodeur, un Vision Transformer (ViT), traite uniquement les patches visibles de l'image (sans tokens de masquage), tandis qu'un décodeur léger reconstruit l'image complète à partir des représentations latentes et des tokens de masquage réintroduits à ce stade. Deuxièmement, un taux de masquage élevé (75\%), bien supérieur aux 15\% utilisés par BERT, justifié par la redondance spatiale des images : un masquage faible permettrait une reconstruction par simple interpolation des patches voisins, sans nécessiter de compréhension globale. La cible de reconstruction est constituée des valeurs de pixels (normalisées par patch) des zones masquées, avec une perte MSE calculée uniquement sur ces zones.

Cette conception asymétrique réduit considérablement le coût de calcul puisque l'encodeur ne traite que 25\% des patches, permettant un gain de vitesse d'entraînement de 3$\times$ et plus, ainsi qu'une réduction de la consommation mémoire facilitant le passage à l'échelle vers des modèles de grande capacité. Les auteurs valident leur approche sur ImageNet-1K en pré-entraînement, puis évaluent par fine-tuning et linear probing. Avec un ViT-Huge, MAE atteint 87.8\% de précision top-1 sur ImageNet-1K (résolution 448), dépassant l'état de l'art parmi les méthodes utilisant uniquement IN1K. Le transfer learning est également démontré sur la détection et la segmentation d'objets sur COCO (53.3 AP\textsubscript{box} avec ViT-L, contre 49.3 pour le pré-entraînement supervisé) ainsi que sur la segmentation sémantique sur ADE20K (53.6 mIoU, contre 49.9 en supervisé).

Les ablations menées par les auteurs apportent plusieurs résultats notables : un encodeur sans tokens de masquage améliore la précision de 14\% en linear probing tout en accélérant l'entraînement, la profondeur du décodeur influence significativement la qualité des représentations latentes (un décodeur suffisamment profond améliore le linear probing de 8\%), et le pré-entraînement bénéficie de schedules longs (1600 epochs sans saturation observée). Les auteurs comparent également différentes stratégies de masquage (aléatoire, par blocs, par grille) ainsi que différentes cibles de reconstruction (pixels, pixels normalisés, tokens dVAE), montrant que les choix simples (masquage aléatoire et pixels normalisés) sont les plus efficaces. MAE se distingue ainsi par sa simplicité conceptuelle, sa scalabilité et son efficacité computationnelle, bien qu'il produise des représentations moins linéairement séparables que les méthodes contrastives, nécessitant un fine-tuning pour être pleinement exploité.

\end{document}